\mychapter{Partie III : Modélisation de Séquences}{Partie III : Modélisation de Séquences}

\section{Contexte et préparation textuelle}
\slidepurpose{Modéliser des données séquentielles avec dépendances temporelles.}{Partie 3 / Section 1}

La dernière partie s'attaque aux données textuelles (NLP) issues de critiques de type IMDb, pour une tâche de traduction automatique de l'anglais vers le français. Les séquences diffèrent fondamentalement des images car leur longueur est variable et l'information dépend de l'ordre temporel des mots. 

La préparation a nécessité la création de vocabulaires distincts (Source et Cible) basés sur la fréquence des mots, incluant des \textit{tokens} spéciaux (\texttt{<unk>}, \texttt{<pad>}, \texttt{<bos>}, \texttt{<eos>}), afin d'encoder les séquences textuelles en identifiants numériques.

\section{Comparaison structurelle : RNN, LSTM, GRU}

L'étude des états cachés a mis en évidence les différences d'architecture entre un simple RNN, une cellule LSTM et une cellule GRU. 

\begin{itemize}
    \item \textbf{RNN :} État simple \texttt{(1, 2, 32)}. Sujet au problème de disparition du gradient.
    \item \textbf{LSTM :} L'état est un tuple \texttt{(hn, cn)}, soit un \textit{hidden state} et un \textit{cell state} (chacun de taille \texttt{(1, 2, 32)}), agissant comme une "bande transporteuse" d'information longue distance.
    \item \textbf{GRU :} Version optimisée combinant les portes d'oubli et de mise à jour, avec un seul état caché \texttt{(1, 2, 32)}, offrant un compromis robuste entre performance et légèreté computationnelle.
\end{itemize}

\section{Architecture Seq2Seq et Entraînement}

Pour la traduction, un modèle encodeur-décodeur (\texttt{Seq2Seq}) basé sur des cellules GRU a été déployé. 

\textbf{Définition du processus de traduction (Extrait Seq2Seq) :}
\begin{verbatim}
class Seq2Seq(nn.Module):
    def forward(self, src, tgt):
        # 1. L'encodeur traite la phrase source en entier
        enc_outputs = self.encoder(src)
        
        # 2. Le decodeur s'initialise avec l'etat final de l'encodeur
        dec_state = self.decoder.init_state(enc_outputs)
        
        # 3. Le decodeur genere la sequence cible
        return self.decoder(tgt, dec_state, enc_outputs)
\end{verbatim}

L'entraînement sur 120 époques intègre le \textbf{Teacher Forcing} (alimentation de la vraie étiquette au temps $t-1$ pour générer la sortie au temps $t$) et le \textbf{Gradient Clipping} (\texttt{clip\_grad\_norm\_}) fixé à 1.0 pour prévenir l'explosion des gradients (problème structurel à la \textit{Backpropagation Through Time}). La perte moyenne finale (Loss) est descendue à 0.0037, garantissant une \textbf{perplexité parfaite de 1.00}.

\section{Stratégies de Décodage et Score BLEU}

Deux stratégies de génération ont été comparées pour la traduction d'une phrase de test : \textit{"this movie is great and fantastic"} dont la vérité terrain est \textit{"ce film est formidable et fantastique"}.

\begin{enumerate}
    \item \textbf{Décodage Glouton (Greedy Decoding)} : Le modèle sélectionne systématiquement le \textit{token} ayant la probabilité maximale au temps courant. Cette approche locale peut manquer l'optimum global.
    \item \textbf{Recherche par faisceau (Beam Search, $k=3$)} : Le modèle explore simultanément les $k$ séquences les plus probables, gardant une trace cumulative des probabilités pour favoriser une génération plus cohérente.
\end{enumerate}

Le \textbf{score BLEU} (Bilingual Evaluation Understudy), reposant sur une analyse des n-grammes avec pénalité de brièveté, a permis d'évaluer quantitativement les prédictions :

\begin{center}
\renewcommand{\arraystretch}{1.3}
\begin{tabular}{l c}
\rowcolor{TableHeader} \tblhead{Stratégie de décodage} & \tblhead{Score BLEU Obtenu} \\
\rowcolor{TableAlt} \textbf{Décodage Glouton} & 0.3333 \\
\textbf{Beam Search (k=3)} & 0.3333 \\
\end{tabular}
\end{center}

\begin{alertblock}{Analyse des prédictions}
Dans cette expérimentation, le score BLEU des deux méthodes est identique. Le décodage glouton a généré \textit{"ce film est formidable et fantastique et fantastique fantastique et"}, démontrant une tendance à la répétition (boucle infinie caractéristique des RNN). Le Beam Search, bien qu'il ait produit \textit{"ce film est formidable et et fantastique et fantastique fantastique"}, souffre du même défaut dans cette configuration, soulignant la complexité d'un arrêt prématuré \texttt{<eos>} correctement inféré.
\end{alertblock}

\solutionfigure{figures/fig_p3_recap.png}{Évaluation des stratégies de décodage (Exemples de métriques typiques NLP)}
